% --- PREÁMBULO FINAL OPTIMIZADO PARA LA EDITORIAL UD ---

% 1. CLASE DEL DOCUMENTO Y OPCIONES GLOBALES (PRIORIDAD DE LA PLANTILLA)
% NOTA: La clase "book" es la base, las opciones son las del formato UD.
%\documentclass[12pt,twoside,spanish]{book}

% 2. CONFIGURACIÓN DEL LENGUAJE Y TIPOGRAFÍA UD (PRIORIDAD AL DISEÑO)

\usepackage{fontspec} % Necesario para LuaLaTeX y la tipografía UD
\input{esConfigDir/esConfigFonts} % Definición de fuentes (lhseries, medseries, etc.)


% 3. TIPOGRAFÍA Y ESTILOS DE LA EDITORIAL UD (FUENTES Y GEOMETRÍA)
% NOTA: Se elimina tu \usepackage{geometry} y tu \usepackage{fontspec} manual
% para usar la configuración exacta de la editorial.


\usepackage{geometry}
\geometry{
	paperheight=24.0cm,
	paperwidth=17.0cm,
	textwidth=12.5cm,
	tmargin=2.5cm,
	bmargin=2.5cm,
	left=2.5cm,
	right=2.0cm
	%   bindingoffset=1.1cm % Usar si se desea ajuste al encuadernar
}

% 4. ESTILOS GRÁFICOS Y COMANDOS UD (\portadillauno, \espart, \nota, \fuente, 'esstyle')

%%% Editar %%%
\def\texttitulocorto{Título corto del libro}
%Lista de autores
\def\loauthors{Nombre o nombres de los autores del libro separados por comas}
%%%%%%%%%%%%%%

% Estos son necesarios para las portadillas y las cornisas especiales.
\input{esConfigDir/esEstilo}

%\usepackage{showframe}% Similar al calderon, apagar al terminar de maquetar.
%\usepackage[cam, width=190truemm, height=260truemm, pdftex,center]{crop} % Descomentar solo en producción final

% 5. AJUSTES FINALES DE FORMATO DEL PÁRRAFO
\usepackage[spanish]{babel}
\usepackage{csquotes}
\usepackage{microtype} % Microtipografía avanzada
\usepackage{ragged2e} % Para control de alineación
\usepackage{authblk} % Manejo de autores
\decimalpoint % Para usar la coma como separador decimal.
\setlength\parindent{0pt} % Sin sangría al inicio de párrafos, como pide la UD.
\newcommand{\normalfonts}{\fontsize{12pt}{14.4pt}\selectfont}
\let\normalsize\normalfonts



%%%%%%%%%%%%%%%%%%%%%%%%%%%%%%%%%%%%%%%%%%%%%%%%%%%%%%%%%%%%%%%%%%%%%%%%%%%%%%%%%%%%%%%%%%%%%%%%%
%%% 6. BIBLIOGRAFÍA (PRIORIDAD DE LA PLANTILLA - USANDO TU ESTILO IEEE)
%%%%%%%%%%%%%%%%%%%%%%%%%%%%%%%%%%%%%%%%%%%%%%%%%%%%%%%%%%%%%%%%%%%%%%%%%%%%%%%%%%%%%%%%%%%%%%%%%
% NOTA: Se usa el backend biber y se mantiene tu estilo IEEE, pero con las correcciones necesarias
% para evitar errores de comandos con la plantilla UD (si se llegara a usar APA).

\usepackage[
backend=biber,
style=ieee, % Se mantiene tu estilo IEEE
]{biblatex}

% Confirma que la ruta de tu archivo .bib es correcta.
\addbibresource{references.bib}
%%%%%%%%%%%%%%%%%%%%%%%%%%%%%%%%%

% CORRECCIÓN IEEE para idioma español (Soluciona el conflicto de la "y" vs "and")
\DefineBibliographyStrings{spanish}{%
	and = {and}%
}
\DeclareDelimFormat{finalnamedelim}{%
	\ifnumgreater{\value{liststop}}{2}{\finalandcomma\addspace}{}%
	\bibstring{and}\addspace}
\DefineBibliographyStrings{english}{%
	andothers = {et\addabbrvspace al\adddot}
}

% ÍNDICES Y SUBARCHIVOS
\usepackage{makeidx}
\usepackage{subfiles}
%%%%%%%%%%%%%%%%%%%%%%%%%%%%%%%%%%%%%%%%%%%%%%%%%%%%%%%%%%%%%%%%%%%%%%%%%%%%%%%%%%%%%%%%%%%%%%%%%
%%% 7. CORNISAS Y ESTILO DE PÁGINA (PRIORIDAD DE LA PLANTILLA Y TUS REQUERIMIENTOS)
%%%%%%%%%%%%%%%%%%%%%%%%%%%%%%%%%%%%%%%%%%%%%%%%%%%%%%%%%%%%%%%%%%%%%%%%%%%%%%%%%%%%%%%%%%%%%%%%%

% NOTA: Se ignora la sección "Cornizas personalizadas" de tu código anterior y se usa el estilo
% de la plantilla, que se activa con \pagestyle{esstyle} en el cuerpo del documento.
\usepackage{fancyhdr}
\usepackage{emptypage}

% Definición del estilo 'plain' para páginas de inicio de capítulo (sin cornisas)
%\fancypagestyle{plain}{
	%\fancyhf{} % Limpiar todo
%\renewcommand{\headrulewidth}{0pt} % Quitar la línea
%\fancyfoot[C]{\thepage} % Número de página centrado en el pie (estándar)
%}
% Definición del estilo 'plain' - AHORA COINCIDE CON EL PIE DE esstyle
\fancypagestyle{plain}{%
	\fancyhf{} % Limpiar todo (incluyendo el número centrado anterior)
\renewcommand{\headrulewidth}{0pt} % Quitar la línea
	% Pie de página izquierdo (Páginas Pares, E.S. | #)
\fancyfoot[LE]{{\fontsize{10.5}{1}\usefont{T1}{Roboto-TLF}{thin}{}\selectfont{E}}\hspace{-0.5mm}{\footnotesize\usefont{T1}{Inter-Sup}{thin}{}\selectfont{\reflectbox{S}}}{\usefont{T1}{Roboto-TLF}{thin}{}\selectfont{\thinspace|\thinspace}}{\lhseries\scriptsize\thepage}}
% Pie de página derecho (Páginas Impares, # | E.S.)
\fancyfoot[RO]{{\lhseries\scriptsize\thepage}{\usefont{T1}{Roboto-TLF}{thin}{}\selectfont{\thinspace|\thinspace}}{\fontsize{10.5}{1}\usefont{T1}{Roboto-TLF}{thin}{}\selectfont{E}}\hspace{-0.5mm}{\footnotesize\usefont{T1}{Inter-Sup}{thin}{}\selectfont{\reflectbox{S}}}}
\fancyfoot[C]{} % Dejar esto vacío para asegurar que el centro está limpio
}



%Esto asegura que la página izquierda de un capítulo nuevo esté vacía (sin cornisas)
\renewcommand{\cleardoublepage}{%
\clearpage\ifodd\value{page}\else\hbox{}\thispagestyle{empty}\newpage\fi%
}

%%%%%%%%%%%%%%%%%%%%%%%%%%%%%%%%%%%%%%%%%%%%%%%%%%%%%%%%%%%%%%%%%%%%%%%%%%%%%%%%%%%%%%%%%%%%%%%%%
%%% 8. PAQUETES FUNCIONALES Y DE CONTENIDO (MANTENIDOS)
%%%%%%%%%%%%%%%%%%%%%%%%%%%%%%%%%%%%%%%%%%%%%%%%%%%%%%%%%%%%%%%%%%%%%%%%%%%%%%%%%%%%%%%%%%%%%%%%%

% PAQUETES DE MATEMÁTICAS Y FÍSICA (TUS PAQUETES ORIGINALES)
\usepackage{amsmath}
\usepackage{amssymb}
\usepackage{amsfonts}
\usepackage{amsthm}
\usepackage{bm}
\usepackage{physics}
\usepackage{fontawesome}
\usepackage{fp}
\usepackage{calculus}
\usepackage{calculator}
\usepackage{fancyvrb}
\usepackage{mathtools}
\usepackage{bigints}
\usepackage{fontawesome}

% PAQUETES DE GRÁFICOS E IMÁGENES (TUS PAQUETES ORIGINALES)
\usepackage{graphicx}
%\usepackage{epstopdf}
%\usepackage{auto-pst-pdf}
\usepackage[export]{adjustbox}
\graphicspath{{./Fig}{./Img}{./Imgn}}

% PAQUETES PSTRICKS
\usepackage{pstricks}
\usepackage{pstricks-add}
\usepackage{pst-node}
\usepackage{pst-3d}
\usepackage{pst-math}
\usepackage{pst-coil}
\usepackage{pst-3dplot}
\usepackage{pst-solides3d}
\usepackage{pdfpages}
% COLOR Y LAYOUT (MANTENIDOS)
\usepackage{xcolor}
\definecolor{orcidlogocol}{HTML}{A6CE39}
\usepackage{academicons}
\usepackage{tcolorbox}
\tcbuselibrary{breakable,skins}
\usepackage{multicol}
\usepackage{esvect}
\usepackage{relsize}
\usepackage{pifont}
\usepackage{etoolbox}
\newtcolorbox{graybox}{
	enhanced,
	boxrule=0pt,
	colback=gray!60!white,
	colupper=white,
	sharp corners,
	boxsep=5pt,
	left=1em, right=0.5em, top=0.5em, bottom=0.5em,
	nobeforeafter,
}

% FLOATS Y CAPTIONS (MANTENIDOS)
\usepackage{float}
\usepackage{placeins}
\usepackage{caption}
\usepackage{wrapfig}
\usepackage{array}
\usepackage{subcaption}


%%%%%%%%%%%%%%%%%%%%%%%%%%%%%%%%%%%%%%%%%%%%%%%%%%%%%%%%%%%%%%%%%%%%%%%%%%%%%%%%%%%%%%%%%%%%%%%%%
%%% 9. CONFIGURACIONES Y COMANDOS (MANTENIDOS Y AJUSTADOS)
%%%%%%%%%%%%%%%%%%%%%%%%%%%%%%%%%%%%%%%%%%%%%%%%%%%%%%%%%%%%%%%%%%%%%%%%%%%%%%%%%%%%%%%%%%%%%%%%%

% CONFIGURACIÓN DE LEYENDAS (CAPTIONS)
\captionsetup[figure]{labelfont=bf, labelsep=period, font=footnotesize, justification=centering, width = 0.8\textwidth, name={Figura}}
\captionsetup[table]{labelfont=bf, labelsep=period, font=footnotesize, justification=centering, width = 0.8\textwidth, name={Tabla}}
\usepackage{enumitem}
\usepackage{listings}
\usepackage{needspace}


% COMANDOS PERSONALIZADOS Y ENTORNOS (TUS DEFINICIONES ORIGINALES)
\newcommand{\ds}{\displaystyle}
\newcommand{\bi}{\int}
\newcommand{\la}{\langle}
\newcommand{\ra}{\rangle}
\newcommand{\ei}[1]{\mathbf{e}_{#1}}
\newcommand{\ptc}{\tcbset{colback=gray!10!white,colframe=black!85!blue}\begin{tcolorbox}[breakable,enhanced,title=\bf Parametrizando el sólido]}
\newcommand{\ptr}{\tcbset{colback=gray!10!white,colframe=black!85!blue}\begin{tcolorbox}[breakable,enhanced,title=\bf Parametrizando la región de Integración]}
\newcommand{\tc}{\begin{tcolorbox}[breakable,enhanced,title=\bf Instrucción en {\tt Mathematica}]}
\newcommand{\etc}{\end{tcolorbox}}
\newcommand{\ve}[1]{\mathbf{#1}}
\DeclareMathOperator{\arcsinh}{arcsinh}
\DeclareMathOperator{\arccosh}{arccosh}
\DeclareMathOperator{\arctanh}{arctanh}
\DeclareMathOperator{\arccoth}{arccoth}
\DeclareMathOperator{\arcsech}{arcsech}
\DeclareMathOperator{\arccsch}{arccsch}
\DeclareMathOperator{\rot}{rot}
\newcommand{\cn}[3]{\pscircle(#1,#2){5pt}\rput(#1,#2){$#3$}}
\newcommand{\fig}{\begin{figure}}
\newcommand{\ef}[1]{\caption{#1}\end{figure}}
\newcommand{\vi}{\mathbf{i}}
\newcommand{\vj}{\mathbf{j}}
\newcommand{\vk}{\mathbf{k}}
\newcommand{\brake}{\\[4pt]} % Se ha corregido la sintaxis de este comando también% Comando auxiliar para saltos
		
% ENTORNOS DE TEOREMA Y REDEFINICIONES
% Alias Global para la Función Seno (sin)
\let\sen\sin
\newtheorem{teorema}{Teorema}[section]
\newtheorem{prop}{Proposición}[section]
\newtheorem{defi}{Definición}[section]
\newtheorem{ejer}{Ejercicio}[section]
\newtheorem{ejemplo}{Ejemplo}[section]
\renewcommand{\proofname}{Demostración}
%\renewcommand{\proof}{Demostración} % Redefinición para Español
%Nueva redefinición al Español
\newcommand{\iniciodemostracion}{\vspace{1em}\textbf{Demostración.}\hspace{1em}}
\newcommand{\findemostracion}{\hfill$\square$}
\renewenvironment{proof}
{\iniciodemostracion\normalfont} 
{\findemostracion} % 

%\renewcommand*{\tablename}{\bf Tabla} % Estilo UD
%\renewcommand*{\figurename}{\bf Figura} % Estilo UD

		
% AJUSTES DE ÍNDICE Y PROFUNDIDAD DE SECCIONES
\counterwithout{footnote}{chapter}
\setcounter{tocdepth}{3}
\setcounter{secnumdepth}{3}
\renewcommand*\contentsname{\lhseries Contenido} % Estilo UD para Contenido
\renewcommand{\cfttoctitlefont}{\hfill\Huge}
\renewcommand{\cftdot}{}
		
% SOLUCIÓN DE CONFLICTOS
\AtBeginDocument{\RenewCommandCopy\qty\SI} % Solución de conflicto physics
\makeindex
		
% REDEFINICIÓN DEL FORMATO Y ESPACIADO DE SECCIONES (TÍTULOS)
% Se mantiene tu formato específico de títulos
\usepackage{titlesec}
\usepackage[hidelinks]{hyperref}

\titlespacing{\section}{0.75cm}{5mm}{2mm}
\titlespacing{\subsection}{0cm}{5mm}{2mm} 
\titlespacing{\subsubsection}{0mm}{5mm}{2mm}

\titleformat{\section}[block]{\color{MSBlack}\fontsize{16pt}{2.0mm} \fontssec}
{\thesection.\quad}{0pt}{\raggedright}

\titleformat{\subsection}{\color{MSBlack}\fontsize{13pt}{2.0mm}\fontssubsec}
{\thesubsection.\quad}{12pt}{}

\titleformat{\subsubsection}[block]{\normalfont\bfseries\fontsize{13pt}{2.0mm}}{}{0pt}{\thesubsubsection\quad}
		
% *** RESTAURACIÓN DE NUMERACIÓN ***
% (Se mantienen tus líneas de restauración de TOC para control granular)
\makeatletter
\renewcommand\l@chapter{\@dottedtocline{0}{0em}{1.5em}} % Nivel 0: Capítulo
\renewcommand\l@section{\@dottedtocline{1}{1.5em}{2.3em}} % Nivel 1: Sección
\renewcommand\l@subsection{\@dottedtocline{2}{3.8em}{3.2em}} % Nivel 2: Subsección
\renewcommand\l@subsubsection{\@dottedtocline{3}{5.8em}{4.1em}} % Nivel 3: Subsubsección
\makeatother

\renewcommand{\thesection}{\thechapter.\arabic{section}}
\renewcommand{\thesubsection}{\thesection.\arabic{subsection}}
\renewcommand{\thesubsubsection}{\thesubsection.\arabic{subsubsection}}
		
% --- FIN DEL PREÁMBULO ---
\counterwithout{footnote}{chapter}